\chapter{Neovim}\label{ch:nvim}

Everything up to this point is Vim, and works identically in Neovim. This
chapter is about what Neovim adds.

\section{What Neovim Actually Changes}

Neovim began in 2014 as a fork of Vim aimed at a refactored, more hackable
codebase. The user-visible consequences are these:

\begin{keys}[0.26]
	\vocab{Lua}              & A real programming language for configuration and plugins, alongside Vimscript \\
	\vocab{Built-in LSP}     & A Language Server Protocol client in core: completion, go-to-definition, rename, diagnostics \\
	\vocab{Treesitter}       & Incremental parsing, giving syntax-aware highlighting, indentation and text objects \\
	\vocab{Terminal emulator} & A real shell inside a buffer                                                  \\
	\vocab{Async jobs}       & Plugins can run external processes without blocking the UI                     \\
	\vocab{Remote UI}        & The editor is decoupled from its front end, via a documented RPC API           \\
	\vocab{Better defaults}  & Sensible options on by default; see \cmd{:help nvim-defaults}                  \\
\end{keys}

\begin{remark}
	Vim has since gained its own versions of some of these (\texttt{+terminal},
	\texttt{+job}, Vim9script). The practical difference today is the ecosystem:
	almost all new plugin development targets Neovim's Lua API.
\end{remark}

\begin{note}[Checking your installation]
	\cmd{:version} prints the build. \cmd{:checkhealth} is more useful --- it runs
	every installed module's self-test and reports missing providers, absent
	external tools, and misconfigured plugins. Run it first whenever something
	behaves strangely.
\end{note}

\section{Configuration in Lua}\label{sec:lua-config}

Neovim reads \file{\textasciitilde/.config/nvim/init.lua} at startup (or
\file{init.vim}, but not both). Any Lua file under
\file{\textasciitilde/.config/nvim/lua/} is importable as a module:
\file{lua/config/options.lua} becomes \texttt{require("config.options")}.

\subsection{Options}

\begin{luacode}
vim.opt.number = true            -- :set number
vim.opt.relativenumber = true    -- :set relativenumber
vim.opt.shiftwidth = 2           -- :set shiftwidth=2
vim.opt.expandtab = true
vim.opt.ignorecase = true
vim.opt.smartcase = true
vim.opt.undofile = true
vim.opt.scrolloff = 8

-- list-like options support append/remove
vim.opt.clipboard:append("unnamedplus")
vim.opt.shortmess:append("c")

-- buffer- and window-local variants
vim.opt_local.conceallevel = 0
\end{luacode}

\begin{note}
	\texttt{vim.opt} wraps options in a table-like object that understands lists
	and sets, which is why \texttt{:append()} works. \texttt{vim.o},
	\texttt{vim.bo} and \texttt{vim.wo} are the plain scalar forms for global,
	buffer-local and window-local options respectively.
\end{note}

\subsection{Key mappings}

\begin{luacode}
-- vim.keymap.set(mode, lhs, rhs, opts)
vim.keymap.set("n", "<leader>w", "<cmd>write<cr>", { desc = "Save file" })
vim.keymap.set("n", "<C-h>", "<C-w>h", { desc = "Go to left window" })
vim.keymap.set("v", "<", "<gv", { desc = "Dedent and keep selection" })
vim.keymap.set("t", "<Esc><Esc>", "<C-\\><C-n>", { desc = "Exit terminal mode" })

-- a Lua function as the right-hand side
vim.keymap.set("n", "<leader>d", function()
  vim.diagnostic.open_float()
end, { desc = "Line diagnostics" })

-- delete a mapping
vim.keymap.del("n", "<leader>w")
\end{luacode}

The \texttt{mode} argument is \texttt{"n"} Normal, \texttt{"i"} Insert,
\texttt{"v"} Visual, \texttt{"x"} Visual-only, \texttt{"t"} Terminal,
\texttt{"c"} Command-line, or a table of several. Useful options are
\texttt{desc} (shown by \texttt{which-key}), \texttt{silent},
\texttt{buffer = 0} to scope the mapping to the current buffer, and
\texttt{remap = true} on the rare occasion you want the right-hand side itself
remapped.

\begin{definition}[Leader]
	\key{<leader>} is a user-chosen prefix key, letting you invent mappings without
	colliding with built-ins. It is \key{<Space>} in most modern configurations
	(LazyVim included), with \key{<localleader>} --- used for filetype-specific
	mappings such as VimTeX's --- set to \key{\textbackslash}. Set them
	\emph{before} loading any plugin:
	\begin{luacode}
vim.g.mapleader = " "
vim.g.maplocalleader = "\\"
	\end{luacode}
\end{definition}

\subsection{Autocommands}

An \vocab{autocommand} runs code when an event fires --- opening a file,
changing filetype, writing, entering Insert mode.

\begin{luacode}
-- Highlight the yanked region briefly
vim.api.nvim_create_autocmd("TextYankPost", {
  callback = function()
    vim.highlight.on_yank({ timeout = 200 })
  end,
})

-- Disable conceal in LaTeX buffers, so raw math delimiters stay visible
vim.api.nvim_create_autocmd("FileType", {
  pattern = "tex",
  callback = function()
    vim.opt_local.conceallevel = 0
  end,
})

-- Groups make a config reloadable without stacking duplicates
local group = vim.api.nvim_create_augroup("MyConfig", { clear = true })
vim.api.nvim_create_autocmd("BufWritePre", {
  group = group,
  pattern = "*.lua",
  callback = function() vim.lsp.buf.format() end,
})
\end{luacode}

\begin{note}
	\cmd{:autocmd} with no argument lists everything currently registered, and
	\cmd{:verbose autocmd BufWritePre} says which file defined each one --- the
	fastest way to find out which plugin is reformatting your code behind your
	back. The same \cmd{:verbose} trick works for options (\cmd{:verbose set
		shiftwidth?}) and mappings (\cmd{:verbose nmap <leader>w}).
\end{note}

\subsection{Talking to Vim from Lua}

\begin{keys}[0.34]
	\texttt{vim.cmd("colorscheme x")} & Run an Ex command                          \\
	\texttt{vim.cmd.colorscheme("x")} & The same, in method form                    \\
	\texttt{vim.fn.expand("\%:p")}    & Call any Vimscript function                 \\
	\texttt{vim.g.foo}                & Global variables (\cmd{g:foo})              \\
	\texttt{vim.b}, \texttt{vim.w}, \texttt{vim.t} & Buffer, window, tab variables  \\
	\texttt{vim.api.nvim\_*}          & The full RPC API --- buffers, windows, extmarks \\
	\texttt{vim.uv}                   & The libuv event loop: timers, filesystem, processes \\
	\texttt{vim.system(\{"ls"\})}     & Run a process asynchronously                \\
	\texttt{vim.notify("msg")}        & Emit a notification                         \\
	\texttt{vim.inspect(tbl)}         & Pretty-print a table                        \\
\end{keys}

\begin{eg}
	At the command line, \cmd{:lua} runs a line of Lua and \cmd{:=} prints an
	expression: \cmd{:=vim.opt.shiftwidth:get()} or
	\cmd{:=vim.api.nvim\_buf\_get\_name(0)}. For anything longer, \cmd{:Lazy} has
	a REPL-like scratch buffer, and \cmd{:source \%} re-runs the Lua file you are
	editing.
\end{eg}

\hrulebar

\section{The Built-in Terminal}\label{sec:terminal}

\begin{keys}[0.30]
	\cmd{:terminal}               & Open a shell in the current window        \\
	\cmd{:split \textbar{} terminal} & \dots{} in a horizontal split          \\
	\key{i} / \key{a}             & Enter Terminal mode (keys go to the shell) \\
	\key{<C-\textbackslash>} \key{<C-n>} & Back to Normal mode                \\
	\cmd{:terminal make}          & Run a command and keep its output in a buffer \\
\end{keys}

A terminal buffer is an ordinary buffer once you are in Normal mode, so every
motion, search and yank applies to the scrollback. That alone justifies using
it over a separate terminal window when you want to copy a stack trace into
your code.

\begin{note}
	LazyVim binds \key{<C-/>} to a floating toggleable terminal, and
	\key{<leader>} \key{f} \key{t} to one rooted at the project directory.
\end{note}

\section{Treesitter}\label{sec:treesitter}

\begin{definition}[Treesitter]
	An incremental parser generator. Neovim keeps a real \emph{syntax tree} of each
	buffer, updated on every keystroke, rather than the regex-based approximation
	that classical Vim syntax files use.
\end{definition}

What that buys you:

\begin{keys}[0.30]
	Highlighting  & Correct on nested and embedded languages --- Lua inside Vimscript, code blocks inside Markdown \\
	Indentation   & Derived from the tree instead of heuristics                                                    \\
	Text objects  & \key{a} \key{f} for a function, \key{i} \key{c} for a class, via \texttt{nvim-treesitter-textobjects} \\
	Incremental selection & Grow the Visual selection node by node up the tree                                      \\
	Folding       & \cmd{'foldexpr'} set from the tree structure                                                    \\
	Navigation    & Jump to the next function or class, and to matching structural pairs                            \\
\end{keys}

\begin{keys}[0.30]
	\cmd{:TSInstall lua}     & Install a parser                                  \\
	\cmd{:TSUpdate}          & Update every installed parser                     \\
	\cmd{:InspectTree}       & Show the live syntax tree for this buffer         \\
	\cmd{:Inspect}           & Report every highlight group under the cursor     \\
\end{keys}

\begin{eg}
	\cmd{:Inspect} is the tool to reach for when a colour looks wrong: it names
	the Treesitter capture, the syntax group and the extmark responsible for the
	character under the cursor, which turns ``my theme is broken'' into a specific
	highlight group you can override.
\end{eg}

\section{LSP: Language Servers}\label{sec:lsp}

\begin{definition}[Language Server Protocol]
	A JSON-RPC protocol, originally from VS Code, that lets one language-analysis
	server (\texttt{lua\_ls}, \texttt{pyright}, \texttt{texlab}, \texttt{clangd})
	serve any editor that speaks it. Neovim ships the client; you supply the
	servers.
\end{definition}

The default mappings once a server attaches to a buffer:

\begin{keys}[0.26]
	\key{K}                & Hover documentation                       \\
	\key{g} \key{d}        & Go to definition                          \\
	\key{g} \key{D}        & Go to declaration                         \\
	\key{g} \key{r}        & List references                           \\
	\key{g} \key{i}        & Go to implementation                      \\
	\key{<C-o>}            & Jump back afterwards                      \\
	\key{<leader>} \key{c} \key{a} & Code action (LazyVim binding)     \\
	\key{<leader>} \key{c} \key{r} & Rename symbol                     \\
	\key{<C-s>} in Insert mode     & Signature help                    \\
\end{keys}

Diagnostics are handled by a separate module, \texttt{vim.diagnostic}, so they
work with linters that are not language servers:

\begin{keys}[0.38]
	\key{]} \key{d} / \key{[} \key{d} & Next / previous diagnostic       \\
	\cmd{:=vim.diagnostic.open\_float()} & Show the full message         \\
	\cmd{:=vim.diagnostic.setloclist()}  & Send them to the location list \\
	\cmd{:LspInfo}         & Which servers are attached to this buffer  \\
	\cmd{:LspLog}          & The client's log, for when they are not    \\
	\cmd{:Mason}           & Install and manage servers, linters, formatters \\
\end{keys}

A minimal hand-rolled setup, for reference --- distributions do this for you:

\begin{luacode}
require("lspconfig").lua_ls.setup({
  on_attach = function(_, bufnr)
    local map = function(keys, fn, desc)
      vim.keymap.set("n", keys, fn, { buffer = bufnr, desc = desc })
    end
    map("gd", vim.lsp.buf.definition, "Goto definition")
    map("K", vim.lsp.buf.hover, "Hover")
    map("<leader>rn", vim.lsp.buf.rename, "Rename")
    map("<leader>ca", vim.lsp.buf.code_action, "Code action")
  end,
  settings = { Lua = { diagnostics = { globals = { "vim" } } } },
})
\end{luacode}

\begin{remark}
	Neovim 0.11 added \texttt{vim.lsp.config()} and \texttt{vim.lsp.enable()},
	which remove the need for \texttt{nvim-lspconfig} as a plugin, and made
	\key{g} \key{r} \key{n} (rename), \key{g} \key{r} \key{a} (code action) and
	\key{g} \key{r} \key{r} (references) default mappings. The machine these notes
	were written on runs 0.9.5, where the \texttt{lspconfig} form above still
	applies --- check \cmd{:version} before copying configuration off the internet,
	because this is the API that has moved fastest.
\end{remark}

\section{Completion and Snippets}

Neovim's built-in completion is \key{<C-n>} / \key{<C-p>} (words in the buffer),
\key{<C-x>} \key{<C-f>} (filenames) and \key{<C-x>} \key{<C-o>} (omni-completion,
which the LSP client provides). Most setups layer a plugin on top ---
\texttt{nvim-cmp} or \texttt{blink.cmp} --- that merges LSP, buffer, path and
snippet sources into one popup, and \texttt{LuaSnip} for expandable templates.
\autoref{ch:setup} shows this configuration as it actually stands on this
machine.

\section{Plugins}\label{sec:plugins}

\begin{definition}[\texttt{lazy.nvim}]
	The current standard plugin manager. Plugins are declared as Lua tables ---
	``specs'' --- and it handles cloning, updating, lockfiles and, as the name
	suggests, deferring a plugin's load until it is needed.
\end{definition}

\begin{luacode}
return {
  -- 1. minimal: just clone it
  { "tpope/vim-surround" },

  -- 2. with options; lazy.nvim calls require("gitsigns").setup(opts)
  { "lewis6991/gitsigns.nvim", opts = { current_line_blame = true } },

  -- 3. lazy-loaded on a filetype, with dependencies and keymaps
  {
    "nvim-telescope/telescope.nvim",
    dependencies = { "nvim-lua/plenary.nvim" },
    cmd = "Telescope",
    keys = {
      { "<leader>ff", "<cmd>Telescope find_files<cr>", desc = "Find files" },
    },
  },

  -- 4. disable something a distribution enabled
  { "folke/trouble.nvim", enabled = false },
}
\end{luacode}

The loading triggers are \texttt{event} (any autocommand event, such as
\texttt{VeryLazy} or \texttt{BufReadPost}), \texttt{ft} (filetype),
\texttt{cmd}, and \texttt{keys}. A plugin declaring none of them loads at
startup.

\begin{keys}[0.26]
	\cmd{:Lazy}        & The plugin dashboard                              \\
	\cmd{:Lazy sync}   & Install, clean and update to match the spec       \\
	\cmd{:Lazy update} & Update, writing \file{lazy-lock.json}             \\
	\cmd{:Lazy restore} & Roll every plugin back to the lockfile           \\
	\cmd{:Lazy profile} & Per-plugin startup timings                       \\
	\cmd{:Lazy log}    & What changed in the last update                   \\
\end{keys}

\begin{note}
	\file{lazy-lock.json} pins every plugin to an exact commit. Commit it to
	version control with your config: \cmd{:Lazy restore} then reproduces a working
	editor exactly, which is what you want after an update breaks something at an
	inconvenient moment.
\end{note}

\begin{moral}
	Plugins are the fun part and the trap. Every one is code that runs on every
	keystroke and one more thing to debug at 2am. Add them one at a time, and only
	after you have felt the absence they fill.
\end{moral}
